\documentclass[border=0pt]{standalone}
% The shared preamble of every docs diagram (docs/diagrams/<name>.tex).
% dvisvgm's driver first, so patterns and opacity survive into the SVG.
\def\pgfsysdriver{pgfsys-dvisvgm.def}
\usepackage[siunitx,RPvoltages]{circuitikz}
\usepackage{xcolor}
% Latin Modern: outline fonts for every glyph a label uses, µ included.
\usepackage[T1]{fontenc}
\usepackage{lmodern}
\usetikzlibrary{arrows.meta,backgrounds,calc,patterns,positioning}
\tikzset{
  % A white card, so diagrams read in mdBook's dark themes too.
  every picture/.append style={show background rectangle,
    background rectangle/.style={fill=white,inner sep=4pt}},
  % Layout geometry: layers told apart by fill and hatching, never colour alone.
  layer a/.style={draw=black,line width=0.6pt,fill=gray!20},
  layer b/.style={draw=black,line width=0.6pt,fill=gray!45},
  layer c/.style={draw=black,line width=0.6pt,fill=white,postaction={pattern=north east lines}},
  layer d/.style={draw=black,line width=0.6pt,fill=white,postaction={pattern=crosshatch dots}},
  marker/.style={draw=black,dashed,line width=0.6pt},
  % A measured distance and its label.
  dim/.style={{Latex[length=2mm]}-{Latex[length=2mm]},line width=0.5pt},
  dim label/.style={font=\small,fill=white,inner sep=1pt},
  % Where a rule fails: thick, dashed and labelled, not only red.
  violation/.style={draw=red!75!black,line width=1.4pt,dash pattern=on 3pt off 2pt},
  violation label/.style={font=\small\bfseries,text=red!75!black},
  caption/.style={font=\small},
}

\begin{document}
\begin{circuitikz}[x=0.18cm,y=0.18cm]
  % Top view of the example (µm): nwell ring 0-40 (3 µm), ntap ring
  % 0.5-39.5 (1 µm), ndiff 10-14 on the pad's net.
  \fill[layer a,even odd rule] (0,0) rectangle (40,40) (3,3) rectangle (37,37);
  \fill[layer d,even odd rule] (0.5,0.5) rectangle (39.5,39.5) (1.5,1.5) rectangle (38.5,38.5);
  \fill[layer c] (10,10) rectangle (14,14);
  \fill[layer b] (11.5,11.5) rectangle (12.5,12.5);
  \node[caption,anchor=north] at (12,9.5) {ndiff};
  \draw[->] (12,17) node[caption,anchor=south] {pad} -- (12,12.6);
  \node[caption,anchor=north] at (20,-0.5) {ntap ring inside nwell ring};
  % Gap from the target to the ring's inner edge (box to box).
  \draw[dim] (14,12) -- node[dim label] {\SI{24.5}{\micro\metre} > \SI{15}{\micro\metre}} (38.5,12);
  \draw[violation] (9.5,9.5) rectangle (14.5,14.5);
  \node[violation label,anchor=west] at (15,7.5) {IO.1};
  % Ring width.
  \draw[violation] (38,20) rectangle (40,30);
  \node[violation label,anchor=west,align=left] at (41,25) {IO.2: ring\\ \SI{1}{\micro\metre} < \SI{2}{\micro\metre}};
\end{circuitikz}
\end{document}
